\scanpage{101}
% Archive foliation: 102. Reused unrelated typescript; printed page 51, not an author's manuscript page number.
\begin{center}-- 51 --\hfill n\textsuperscript{o} 380\end{center}
3. \uline{Sous-algèbres de} $\mathcal B(\Omega)$.

Soient $\Omega$ un espace topologique et $R$ une relation d'équivalence dans $\Omega$. On notera $T'(R)$ la sous-algèbre fermée involutive de $\mathcal B(\Omega)$ constituée par les fonctions continues bornées $f$ sur \qquad qui sont constantes sur les classes de $R$. Cette sous-algèbre $T'(R)$ est l'image de l'algèbre $\mathcal B(\Omega/R)$ par l'isométrie $q^{\mathcal B}:\mathcal B(\Omega/R)\,\mathcal B(\Omega)$ que définit l'application quotient $q:\Omega\longrightarrow\Omega/R$. D'autre part, pour toute sous-algèbre $A$ de $\mathcal B(\Omega)$, on notera $T(A)$ la relation d'équivalence dans $\Omega$ constituée par les couples $(t,t')\in\Omega^2$ tels que $f(t)=f(t')$ pour tout $f\in A$. Il est clair que le graphe de $T(A)$ est fermé. La relation d'équivalence $T(A)$ n'est pas nécessairement fermée. On a cependant le Lemme suivant :

LEMME 2.-- \uline{Pour toute sous-algèbre} $A\subset\mathcal B(\Omega)$, \uline{le saturé d'une partie compacte de} $\Omega$ \uline{pour la relation d'équivalence} $T(A)$ \uline{est fermé}.

Soit en effet $t$ un point de $\Omega$ adhérent au saturé d'un compact $K\subset\Omega$. Pour toute fonction $f\in A$ et tout $\varepsilon>0$, le voisinage fermé saturé de $f$ constitué par les points $t'\in\Omega$ tels que $|f(t)-f(t')|\leq\varepsilon$ rencontre $K$. Puisque $K$ est compact, l'intersection de ces voisinages (lorsque $f$ parcourt $A$ et que $\varepsilon$ parcourt l'ensemble des nombres réels $>0$) rencontre $K$. Cette intersection étant la classe d'équivalence de $t$, ceci prouve que $t$ appartient au saturé de $K$.

Il est évident que
\[
T'(T(A))\supset A
\]
\[
T(T'(R))\supset R
\]
pour toute sous-algèbre $A$ de $\mathcal B(\Omega)$ et toute relation d'équivalence $R$ sur $\Omega$.

PROPOSITION 4.-- \uline{Si} $\Omega$ \uline{est normal et si} $R$ \uline{est une relation d'équivalence fermée sur} $\Omega$, \uline{alors} $T(T'(R))=R$.

En effet, on sait que $\Omega/R$ est dans ce cas normal (Top. Gén., Chap.9, \S\ 4, exer.15).
